\documentclass[12pt,a4paper]{article}
\usepackage{fontawesome5}
\providecommand{\faPencilAlt}{\faIcon{pencil-alt}}
\usepackage{styles/couleurs}
\usepackage{styles/mise_en_page}
\usepackage{styles/boites}
\usepackage{styles/maths}
\usepackage{styles/figures}
\usepackage{styles/evaluations}
\usepackage{enumitem}
\usepackage{array}
\usepackage{tabularx}
\usepackage{amsmath}
\setlist[itemize]{leftmargin=1.4em,itemsep=0.15em,topsep=0.2em}
\setlist[enumerate]{leftmargin=*,itemsep=0.25em,topsep=0.25em}
\begin{document}

\enteteinterro[Calculatrice interdite]{4e}{Écrire et simplifier}{12 min}{10}

\exo{3}[Représenter]
Traduis par une expression littérale. \pts{3}

\begin{enumerate}[label=\alph*)]
  \item Le triple de \(x\), augmenté de \(5\).
  \item Le produit de \(4\) par la somme de \(x\) et de \(2\).
  \item La différence entre le carré de \(x\) et \(7\).
\end{enumerate}

\lignes{4}

\exo{3}[Calculer]
Pour \(x=-2\), calcule la valeur de : \pts{3}

\[
A=3x+7
\qquad
B=x^2-5
\qquad
C=2(x+4).
\]

\lignes{5}

\exo{3}[Calculer]
Simplifie les écritures. \pts{3}

\[
4\times x
\qquad
x\times7
\qquad
x\times x
\qquad
1\times x
\qquad
3\times(x+2).
\]

\lignes{3}

\exo{1}[Raisonner]
Explique pourquoi \(3+x\) ne peut pas être simplifié en \(3x\). \pts{1}

\lignes{2}

\end{document}
